\subsection{A class-number certificate for the quintic residual field}

The next calculation proves three-primary class-group vanishing
directly and supplies a second proof of source dominance.
It uses the small-prime term in the unconditional explicit formula.

\begin{proposition}\label{prop:quintic-three-class}
For $K=\Q(\Jac(v^2+v=u^5)[3])$, one has
$\operatorname{Cl}(K)[3]=0$.
\end{proposition}
\begin{proof}
If the three-primary class group were nonzero, class field theory
would give an everywhere-unramified cyclic cubic extension $L/K$.
It would be totally complex of degree $240$, with
\[
 \operatorname{rd}(L)=\operatorname{rd}(K)=3^{76/80}5^{94/80}.
\]
The unique prime $\mathfrak p\mid3$ of $K$ has norm $81$, and
$(3)=\mathfrak p^{20}$. Its ideal class has order dividing twenty,
so its Artin image in $\operatorname{Gal}(L/K)$ is trivial.
Consequently $L$ has three distinct primes of norm $81$.

Define an even function $H$, zero outside $[-2,2]$, by
\[
 H(t)=\frac{(2-t)(1+\tfrac12\cos\pi t)}3
       +\frac{\sin\pi t}{2\pi}\quad(0\le t\le2).
\]
For $g(x)=1+\cos(\pi x)$ on $[-1,1]$ and $g(x)=0$ elsewhere,
direct integration gives
\[
 H(t)=\frac13\int_{\mathbb R}g(x)g(x-t)\,dx.
\]
Thus $H\ge0$, $H(0)=1$, $H\le1$ by Cauchy--Schwarz, and its
Fourier transform is nonnegative. Set $b=51/5$ and
\[
 F(x)=\frac{H(x/b)}{\cosh(x/2)},\quad
 I=\int_0^\infty\frac{1-F(x)}{2\sinh(x/2)}\,dx,
 \quad \log B=\log(8\pi)+\gamma-I.
\]
The explicit formula and unconditional positivity criterion of
\citet[Eqs.~(2.3)--(2.4), p.~122]{Odlyzko1990} give, for any
totally complex field of degree $d$,
\[
 \log\operatorname{rd}\ge\log B+\frac{f-E}{d},\qquad
 E=4\int_0^\infty F(x)\cosh(x/2)\,dx,
\]
where
\[
 f=2\sum_{\mathfrak l}\sum_{m\ge1}
 \frac{\log N\mathfrak l}{(N\mathfrak l)^{m/2}}
 F(m\log N\mathfrak l).
\]
The autocorrelation verifies the positivity condition throughout
the critical strip; no hypothesis on the zeros is imposed.
Since $\int_0^2H(t)\,dt=2/3$, we have $E=8b/3=136/5$.
Every prime contribution is nonnegative. Retaining only the first
powers of the three known norm-$81$ primes gives
\[
 f\ge f_{81}:=
 \frac{6\log81}{9\cosh(\log81/2)}H(\log81/b).
\]

Here are the complete endpoint and tail specifications of the
certified integral. Take $\epsilon=1/16384$. Integrate the displayed
meromorphic expression for $(1-F(x))/(2\sinh(x/2))$ on
$[\epsilon,2b]$ by complex-ball quadrature at $192$ bits, with
absolute and relative tolerances $2^{-120}$, using the certified
Arb integration algorithm described by \citet{Johansson2018}.
Beyond $2b$, its integral is exactly $-\log\tanh(b/2)$.
At the origin, $H'(0)=0$ and
$|H''(t)|\le\|g'\|_2^2/3=\pi^2/3<6$.
For $0\le x\le\epsilon$ this gives
\[
 1-H(x/b)\le3x^2/b^2,
 \qquad 1-\operatorname{sech}(x/2)\le x^2/8.
\]
As $b\ge2$ and $2\sinh(x/2)\ge x$, the integrand lies between
zero and $x$. Its omitted endpoint integral is therefore at most
$\epsilon^2/2$. Subtracting this upper bound from the expression
for $\log B$ gives a certified lower bound $\log B_-$.

Arb ball arithmetic, with exact rational inputs and thresholds,
certifies
\begin{gather*}
 B_->\frac{21046}{1000},\qquad f_{81}>\frac{472}{1000},\\
 \exp\!\left(\log B_-+\frac{f_{81}-136/5}{240}\right)
       >\frac{18828}{1000},\qquad
 3^{76/80}5^{94/80}<\frac{18818}{1000}.
\end{gather*}
For orientation, the two last enclosures are respectively
$18.8285116161758712974\ldots$ and
$18.8171976695261105380\ldots$; the proof uses the rational
inequalities. The explicit-formula lower bound for $L$ exceeds
its established root discriminant, a contradiction.
\end{proof}

\begin{corollary}[A ray-class proof of source dominance]
\label{cor:quintic-ray-source}
Proposition~\ref{prop:quintic-three-class} implies the two
surjections in Theorem~\ref{thm:lm-fifteen-source}.
\end{corollary}
\begin{proof}
By the Frattini criterion it is enough to exclude a nontrivial cyclic
cubic extension of $K$, unramified outside $3,5$, and split at its
prime above three. Such an extension is unramified at three and
has at most tame conductor exponent one at the unique prime
$\mathfrak q\mid5$. Its group is a quotient of the ray class group
modulo $\mathfrak q$. The residue group is $\mathbb F_{25}^{\times}$,
whose three-primary part has order three. The global unit $\zeta_3$
reduces to an element of that exact order, so the ray-class exact
sequence and Proposition~\ref{prop:quintic-three-class} leave no
three-primary quotient. The local decomposition at three already
surjects onto $\operatorname{Gal}(K/\Q)$ by
Proposition~\ref{prop:lm-a5-residual}; the full source surjection follows.
\end{proof}
